% ================================================================
% INTRODUCTION
% ================================================================

% ------------------------------------------------------------------
% 1. MOTIVATION
% ------------------------------------------------------------------

% TODO: Opening paragraph — carbon pricing policies (CPP) are one of the main
% tools to reduce GHG emissions. 75+ initiatives globally (World Bank 2023).
% Theoretical motivation is well-understood since Pigou: emissions are an
% externality, first-best requires pricing all emissions.

% TODO: But in practice, CPP are incomplete. Policymakers target a subset of
% polluting sectors and, within sectors, a subset of firms.
% Present Figure 1 here:
%   - Panel A: share of CPP instruments covering each sector (output/figures/panel_a_sector_targeting.pdf)
%   - Panel C: within-sector coverage heterogeneity in the EU ETS (output/figures/panel_c_within_sector_coverage.pdf)

% \begin{figure}[t]
%     \centering
%     \begin{subfigure}[b]{0.48\textwidth}
%         \centering
%         \includegraphics[width=\textwidth]{../output/figures/panel_a_sector_targeting.pdf}
%         \caption{Sectoral coverage across CPP instruments}
%     \end{subfigure}
%     \hfill
%     \begin{subfigure}[b]{0.48\textwidth}
%         \centering
%         \includegraphics[width=\textwidth]{../output/figures/panel_c_within_sector_coverage.pdf}
%         \caption{Within-sector coverage in the EU ETS}
%     \end{subfigure}
%     \caption{Carbon pricing policies are incomplete}
%     \label{fig:cpp_incomplete}
% \end{figure}

% TODO: Cross-country heterogeneity figure/table. Different countries target
% very different sectors and at very different rates.
% [PLACEHOLDER — need data/source for this]

% TODO: Firms operate in production networks. The effect of taxing a subset of
% firms depends not only on who is taxed but on how they are connected to the
% rest of the economy through input-output linkages.

% ------------------------------------------------------------------
% 2. LITERATURE GAP
% ------------------------------------------------------------------

% TODO: The existing literature has studied CPP effects along two dimensions
% that this paper bridges:
%
% (a) Firm-level effects of CPP:
%     - Colmer et al (2025): firm-level effects of EU ETS in France
%     - Dechezlepetre et al (2023): firm-level effects across European countries
%     - These papers estimate how individual firms respond but do not aggregate
%       effects accounting for network interactions.
%
% (b) Aggregate effects of CPP with fixed targeting:
%     - Metcalf and Stock (2023): aggregate emissions reductions from carbon taxes
%     - Kaenzig (2025): macroeconomic effects of carbon pricing
%     - These papers study aggregate outcomes but take the targeting scheme as
%       given and do not explore how different targeting interacts with the I-O
%       structure.
%
% Neither strand takes seriously both (i) the production network and (ii) the
% fact that targeting is a choice variable.

% ------------------------------------------------------------------
% 3. POLICY RELEVANCE
% ------------------------------------------------------------------

% TODO: There are active discussions to expand the scope of existing CPP:
%   - EU ETS 2 (extending to buildings and transport)
%   - CBAM (Carbon Border Adjustment Mechanism)
%   - Expansion of national carbon taxes
% It is therefore natural to ask: what are the aggregate effects of different
% targeting schemes, and which firms should be prioritized?

% ------------------------------------------------------------------
% 4. THIS PAPER'S CONTRIBUTION
% ------------------------------------------------------------------

% TODO: This paper takes a first step at answering the question above, taking
% seriously that (a) firms operate in production networks and (b) CPP are
% incomplete in practice.
%
% We make three contributions:
%
% First, we develop a theory of how targeted CPP affect aggregate emissions in
% production networks. We decompose the aggregate effect into an abatement
% channel and a reallocation channel, and derive a sufficient statistic — the
% covariance between firms' policy exposure and their network-adjusted emission
% intensity — that governs the reallocation channel.
%
% Second, we measure this sufficient statistic using Belgian micro-data,
% combining firm-level emissions data with business-to-business transaction
% records to construct the economy's I-O matrix at the firm level.
%
% Third, we evaluate counterfactual targeting schemes: what would happen if
% Belgium targeted all industrial emissions? What if it targeted the same
% fraction of emissions but chose firms with higher network centrality?

% ------------------------------------------------------------------
% 5. PREVIEW OF RESULTS
% ------------------------------------------------------------------

% TODO: Preview main findings:
%   - The Belgian I-O structure amplifies the abatement effects of the current
%     EU ETS targeting.
%   - A small number of firms contribute disproportionately to this amplification.
%   - Targeting firms with higher network centrality can achieve larger emissions
%     reductions for the same coverage.
%   - [Counterfactual results to be added]

% ------------------------------------------------------------------
% 6. RELATED LITERATURE
% ------------------------------------------------------------------

% TODO: Position the paper relative to:
%   - Production networks: Acemoglu et al (2012), Baqaee & Farhi (2019, 2020),
%     Carvalho & Tahbaz-Salehi (2019)
%   - Environmental economics & carbon pricing: Nordhaus (1994), Metcalf & Stock
%     (2023), Kaenzig (2025), Colmer et al (2025)
%   - Carbon leakage and border adjustments: Fowlie et al (2016), Copeland et al
%     (2022)
%   - Emissions in production networks: Shapiro (2021), Shapiro & Walker (2018),
%     Levinson (2009)
%   - Belgian/European firm-level data: Dhyne et al (2020, 2023), Duprez et al (2023)
